% =====================================================================================
%  Multi-objective assignment results. All numbers from code/phase4/results/*_fullv3.*
%  1,225 runs: 13 policies, 5 batches, 30 seeded repeats, 150 generations, population 60.
% =====================================================================================

\subsection{Multi-objective assignment}
\label{sec:assignment-results}

The protocol runs five operational baselines, standard NSGA-II, NSGA-III as a many-objective
control, the proposed PAA-NSGA-II, its three single-component ablations, and two post-hoc
exploratory combinations. Five arrival batches, 30 seeded repeats each, 150 generations,
population 60. All \num{1225} runs completed with \textbf{zero capacity violations}, so the repair
operator satisfies constraint C2 throughout.

\begin{table}[htbp]
\centering
\caption{Assignment policies on BPI Challenge 2014. Hypervolume is computed against a
reference point fixed as the nadir across all variants and all runs within a batch. Objectives
$f_1$ and $f_4$ are maximised, $f_2$ and $f_3$ minimised. Jain's index is reported alongside
the $f_3$ standard deviation and is not substituted for it.}
\label{tab:assignment}
\small
%% GENERATED by code/figures/make_tables.py. Do not edit by hand.
%% The caption, label and float live in the manuscript; only the numbers are here.
\begin{tabular}{lrrrrr}
\toprule
Policy & Hypervolume & $f_1$ FTR & $f_2$ breach & $f_3$ workload SD & $f_4$ expertise \\
\midrule
Random                         & -- & 0.4890 & 0.4351 & 0.7807 & 0.3562 \\
Round robin                    & -- & 0.4592 & 0.4510 & 1.1796 & 0.3449 \\
FIFO                           & -- & 0.4500 & 0.5523 & 2.5554 & 0.4039 \\
Least loaded                   & -- & 0.4473 & 0.4528 & 0.4984 & 0.3876 \\
Greedy best outcome            & -- & 0.6998 & 0.2885 & 1.6501 & 0.7162 \\
Standard NSGA-II               & 0.0695 & 0.6169 & 0.3146 & 1.0352 & 0.6988 \\
NSGA-III                       & 0.0824 & 0.6455 & 0.2792 & 0.9978 & 0.7556 \\
PAA-NSGA-II                    & 0.0762 & 0.6306 & 0.3014 & 1.2552 & 0.7478 \\
PAA without priority init      & 0.0658 & 0.6156 & 0.3076 & 1.0467 & 0.6644 \\
PAA without adaptive crossover & 0.0767 & 0.6325 & 0.3028 & 1.2518 & 0.7594 \\
PAA without SLA mutation       & 0.0767 & 0.6303 & 0.2964 & 1.2414 & 0.7373 \\
PAA with adaptive mutation     & 0.0766 & 0.6337 & 0.3046 & 1.2710 & 0.7657 \\
NSGA-III + priority init       & 0.0850 & 0.6595 & 0.2834 & 1.1679 & 0.7931 \\
\bottomrule
\end{tabular}



\vspace{0.4em}
{\footnotesize \textsuperscript{*}Post-hoc exploratory combination, designed after observing
the ablation. Reported as exploratory and not claimed as a contribution.}
\end{table}

\paragraph{The proposed operators improve on their own base algorithm.}
Against standard NSGA-II over 150 paired runs, PAA-NSGA-II is better on four of five metrics
by paired Wilcoxon signed-rank with Holm-Bonferroni correction: hypervolume
($+\num{0.00666}$, $p = \num{1.5e-19}$), first-time-right ($+\num{0.01367}$,
$p = \num{1.3e-5}$), breach ($-\num{0.01322}$, $p = \num{1.5e-4}$) and expertise match
($+\num{0.04898}$, $p = \num{9.9e-7}$).

\paragraph{The base algorithm matters more than the modifications.}
NSGA-III, included only as a many-objective control, reaches hypervolume \num{0.0824} against
\num{0.0762} for the proposed method, a paired difference of $\num{-0.00625}$ at
$p = \num{1.5e-23}$. It does so with a front of about \num{18} solutions against \num{60}, so it
attains higher hypervolume from under a third as many points. With five objectives this is the
regime for which reference-point selection was introduced
\citep{debjain2014evolutionary,zheng2024runtime}, and the result is consistent with that
literature rather than in tension with it.

\paragraph{One of three components does the work.}
Table~\ref{tab:ablation} isolates each modification. Removing priority-aware initialisation
costs \num{0.0104} hypervolume, so that component is doing real work. The other two are
not: removing adaptive crossover \emph{raises} hypervolume by \num{0.0005} and removing
SLA-aware mutation raises it by \num{0.0006}, both small enough to sit inside run-to-run
noise but neither of them a contribution. Reporting the three as a bundle would have credited
them equally, and that would have been wrong in both directions.

\begin{table}[htbp]
\centering
\caption{Ablation. Change in hypervolume when a single component is removed from the proposed
method. A negative delta means the component contributes.}
\label{tab:ablation}
%% GENERATED by code/figures/make_tables.py. Do not edit by hand.
%% The caption, label and float live in the manuscript; only the numbers are here.
\begin{tabular}{lrr}
\toprule
Component removed & Hypervolume without & $\Delta$ \\
\midrule
Priority-aware initialisation  & 0.0658 & $-0.0104$ \\
Adaptive crossover             & 0.0767 & $+0.0005$ \\
SLA-aware mutation             & 0.0767 & $+0.0006$ \\
\midrule
None (full method)             & 0.0762 & \\
\bottomrule
\end{tabular}


\end{table}

\begin{figure}[htbp]
\centering
\includegraphics[width=0.72\textwidth]{fig08_ablation.pdf}
\caption{Component ablation. Removing priority-aware initialisation costs \num{0.0104}
hypervolume; removing either of the other two changes it by less than \num{0.001} and in the
opposite direction, which is inside run-to-run noise. Reporting the three as a bundle
would have credited them equally.}
\label{fig:ablation}
\end{figure}

\paragraph{The best configuration is not a significant improvement.}
Placing the one modification that contributes onto the base algorithm that wins gives the
highest hypervolume observed, \num{0.0850}. Against plain NSGA-III the paired difference is
$\num{+0.00254}$ at $p = \num{6.8e-11}$, which does clear $\alpha = \num{0.05}$. It is still
reported as exploratory and is still not claimed, because it was designed after reading the
ablation table on this same data. A post-hoc combination that clears a threshold on the
data that suggested it has not been tested; it has been fitted. Claiming it would be the
same error this paper is about, committed one level up.

\paragraph{A real trade-off, visible in two independent metrics.}
The proposed method loses significantly on workload spread ($+\num{0.22006}$,
$p = \num{1.6e-21}$). Jain's fairness index agrees: \num{0.1207} against
\num{0.1716} for standard NSGA-II. The two metrics were computed independently and could have disagreed. The
method concentrates work on capable groups to gain outcome quality, which is a trade-off an
operator may or may not want. The exponent $p$ in Equation~\ref{eq:powermean} is the parameter
through which that preference is expressed rather than assumed.

\begin{figure}[htbp]
\centering
\includegraphics[width=0.78\textwidth]{fig07_hypervolume_by_variant.pdf}
\caption{Search quality across algorithm variants. The proposed method improves on its own
base algorithm and is in turn outperformed by an off-the-shelf many-objective control.}
\label{fig:hypervolume}
\end{figure}
